\documentclass[10pt]{article}
\usepackage[letterpaper,margin=0.75in]{geometry}
\usepackage[T1]{fontenc}
\usepackage{lmodern}
\usepackage{microtype}
\usepackage{enumitem}
\usepackage{titlesec}
\usepackage[hidelinks]{hyperref}

\setlist{nosep,leftmargin=1.2em}
\titlespacing*{\section}{0pt}{5pt}{2pt}
\titlespacing*{\subsection}{0pt}{4pt}{2pt}
\titleformat{\section}{\normalsize\bfseries}{\thesection}{0.6em}{}
\titleformat{\subsection}{\normalsize\bfseries}{\thesubsection}{0.6em}{}
\setlength{\parskip}{1pt}
\setlength{\parindent}{0pt}

\newcommand{\limn}{\textit{Limitation:}~}

\begin{document}

\begin{center}
{\large\bfseries Does the Final Answer Measure the Agent? Outcome-Only vs.\ Trajectory-Aware\\[1pt] vs.\ Human Evaluation of a Research-Assistant Agent}\\[4pt]
CS 6180 Foundations of Generative AI (Fall 2026) --- Project Proposal\\[2pt]
\textbf{Team member:} \textit{[Student Name]} (individual project)
\end{center}

\section{Introduction}

\textbf{Problem.} LLM agents that search, read and reason over several steps are usually scored only on whether the final answer matches a reference. This project measures \emph{how much} that outcome-only score disagrees with trajectory-aware metrics (which score the search-and-read process) and with human judgment of the same runs, and whether the disagreement changes which model or configuration is judged ``best''.

\textbf{Why it matters.} A correct answer can come from a poor process: the model may recall it from pretraining, cite documents that do not support it, or succeed by luck on one run. A wrong answer can come from a sound search. If outcome-only evaluation hides these cases, leaderboards and outcome-only training rewards~\cite{jin2025searchr1} optimize the wrong signal. Trajectory evaluation could make agents more \emph{reliable}, \emph{safer} and more \emph{interpretable}, but its LLM judges have their own biases and costs.

\textbf{Broader context.} The work sits in \emph{evaluation and benchmarking of generative-AI agents}, overlapping \emph{alignment/safety} (grounded answers) and \emph{LLM-as-a-judge validity}. Recent work argues that outcome-only evaluation is insufficient~\cite{yehudai2025survey,chen2026trace}, but not \emph{how often} it misleads or \emph{for which decisions}.

\textbf{Planned study.} A ReAct~\cite{yao2023react} research-assistant agent searches a frozen HotpotQA~\cite{yang2018hotpotqa} corpus with gold evidence labels, using Qwen3 models (1.7B, 4B, 8B; thinking on/off) on free Kaggle GPUs, over 400 questions $\times$ 3 seeds plus closed-book and single-retrieval baselines. Each trajectory is scored (A) outcome-only, (B) trajectory-aware (evidence recall, citation grounding, redundant calls, LLM step judge), and (C) by blinded human scoring of 150 trajectories with reliability checks.

\textbf{Impact.} \emph{Scientific:} measured disagreement rates with confidence intervals. \emph{Practical:} which cheap signal best predicts human judgment per unit of compute. \emph{Community:} public trajectories, judge outputs and human labels.

\section{Background / Literature Review}

\textbf{AI Agents That Matter~\cite{kapoor2024agents} and HAL~\cite{kapoor2025hal}.} Propose cost-controlled evaluation with accuracy--cost Pareto frontiers; HAL's 21,730 rollouts show more reasoning effort often \emph{reduces} accuracy, and its log inspection reveals hidden behaviors. Strong evaluation hygiene. \limn quality is still task success; logs are read qualitatively, not scored.

\textbf{TRACE~\cite{chen2026trace}.} The closest work: names the ``high-score illusion'' of Pass@1 and scores trajectories by efficiency and evidence grounding, exposing trade-offs accuracy hides. \limn its evaluator is not validated against humans, checked for bias, or tested across repeated runs.

\textbf{AgentProcessBench~\cite{fan2026agentprocessbench}.} 8,509 human step labels (+1/0/$-$1) on 1,000 tool-use trajectories; weak policies look better because they stop early. I reuse its ternary labels. \limn it evaluates judges, not outcome--process disagreement on the same runs.

\textbf{TRAIL~\cite{deshpande2025trail} and MAST~\cite{cemri2025mast}.} Rigorous error taxonomies for agent traces (MAST: 14 failure modes, $\kappa=0.88$). \limn annotation is costly (about 110 min/trace in TRAIL; six experts in MAST) and unrelated to outcome scores.

\textbf{Judging LLM-as-a-Judge~\cite{zheng2023judging}.} LLM judges match human agreement levels but show position, verbosity and self-preference bias; position swapping mitigates this. \limn tested on single-turn chat, not on agent trajectories, which LLM judges now routinely score~\cite{zhuge2024agentjudge,lu2025agentrewardbench}.

\textbf{$\tau$-bench~\cite{yao2024taubench}.} pass$^k$ (all $k$ runs succeed) shows GPT-4o falling from over 60\% to under 25\% at $k=8$; a simple, revealing metric. \limn still outcome-only: it shows \emph{that} runs diverge, not \emph{why}.

\textbf{Search-R1~\cite{jin2025searchr1}.} Trains search agents with RL using only exact-match reward on a reproducible local retriever, which I reuse. \limn this reward cannot tell retrieved answers from memorized ones.

\textbf{Gap.} No study applies outcome, trajectory and human evaluation to the \emph{same} trajectories and measures how often they disagree, whether rankings change at matched cost, and how reliable each evaluator is.

\section{Research Question}

\textbf{Main question:} \emph{On identical agent trajectories, how much do outcome-only and trajectory-aware evaluations disagree, does it change which configuration is judged best, and which automatic evaluator best tracks human judgment per unit of cost?} Hypotheses are pre-registered before full runs:
\begin{itemize}
  \item \textbf{RQ1 (disagreement).} How often is the answer right but the process bad, or vice versa? \emph{H1:} at least 10\% of correct answers are ungrounded, mostly explained by memorization (compared against closed-book accuracy), and less often for larger models. \emph{Measure:} 2$\times$2 table of correctness vs.\ grounding.
  \item \textbf{RQ2 (rankings).} Do the three evaluations rank the 6 configurations alike? \emph{H2:} outcome and human rankings differ at matched token cost (Kendall $\tau<0.8$)~\cite{kapoor2024agents}.
  \item \textbf{RQ3 (validity).} Which automatic metric best predicts human scores, at what cost? \emph{H3:} rule-based metrics (evidence recall, redundant calls) match an LLM step judge in Spearman $\rho$; LLM judges show measurable bias~\cite{zheng2023judging}.
  \item \textbf{RQ4 (reliability).} \emph{H4:} pass$^3$ is below 70\% of pass$^1$ for every configuration~\cite{yao2024taubench}; seed-to-seed trajectory diffs show why.
\end{itemize}
The minimum deliverable is RQ1 and RQ4 on the 4B model.

\textbf{Why worth answering / gap filled.} It turns the accepted claim ``final-answer accuracy is not enough'' into effect sizes with confidence intervals and decision-level consequences, the measurement missing in Section~2. The design keeps every hypothesis falsifiable: paired questions, 3 seeds, a closed-book contamination control, and a human baseline with measured reliability (intra-rater re-test, external second rater for Cohen's $\kappa$). \textbf{Who benefits:} benchmark designers, practitioners choosing models from leaderboards, researchers training agents with outcome rewards, and builders of LLM judges.

\clearpage
\renewcommand{\refname}{\normalsize References}
\begin{thebibliography}{99}\small\setlength{\itemsep}{0pt}
\bibitem{yehudai2025survey} A.~Yehudai, L.~Eden, A.~Li, G.~Uziel, Y.~Zhao, R.~Bar-Haim, A.~Cohan, and M.~Shmueli-Scheuer. Survey on evaluation of LLM-based agents. \emph{Findings of ACL}, 2026. arXiv:2503.16416.
\bibitem{kapoor2024agents} S.~Kapoor, B.~Stroebl, Z.~S.~Siegel, N.~Nadgir, and A.~Narayanan. AI agents that matter. \emph{Transactions on Machine Learning Research}, 2025. arXiv:2407.01502.
\bibitem{kapoor2025hal} S.~Kapoor, B.~Stroebl, P.~Kirgis, N.~Nadgir, Z.~S.~Siegel, et al. Holistic Agent Leaderboard: The missing infrastructure for AI agent evaluation. \emph{ICLR}, 2026. arXiv:2510.11977.
\bibitem{chen2026trace} Y.~Chen, J.~Jiang, J.~Liu, Y.~Zhang, X.~Guo, and I.~King. TRACE: Trajectory-aware comprehensive evaluation for deep research agents. arXiv:2602.21230, 2026.
\bibitem{fan2026agentprocessbench} S.~Fan, X.~Ye, Y.~Huo, Z.-Y.~Chen, Y.~Guo, S.~Yang, W.~Yang, S.~Ye, et al. AgentProcessBench: Diagnosing step-level process quality in tool-using agents. arXiv:2603.14465, 2026.
\bibitem{deshpande2025trail} D.~Deshpande, V.~Gangal, H.~Mehta, J.~Krishnan, A.~Kannappan, and R.~Qian. TRAIL: Trace reasoning and agentic issue localization. arXiv:2505.08638, 2025.
\bibitem{cemri2025mast} M.~Cemri, M.~Z.~Pan, S.~Yang, L.~A.~Agrawal, B.~Chopra, R.~Tiwari, K.~Keutzer, A.~Parameswaran, et al. Why do multi-agent LLM systems fail? \emph{NeurIPS}, 2025. arXiv:2503.13657.
\bibitem{zheng2023judging} L.~Zheng, W.-L.~Chiang, Y.~Sheng, S.~Zhuang, Z.~Wu, Y.~Zhuang, Z.~Lin, Z.~Li, et al. Judging LLM-as-a-judge with MT-Bench and Chatbot Arena. \emph{NeurIPS Datasets and Benchmarks}, 2023. arXiv:2306.05685.
\bibitem{zhuge2024agentjudge} M.~Zhuge, C.~Zhao, D.~Ashley, W.~Wang, D.~Khizbullin, Y.~Xiong, Z.~Liu, E.~Chang, et al. Agent-as-a-Judge: Evaluate agents with agents. arXiv:2410.10934, 2024.
\bibitem{lu2025agentrewardbench} X.~H.~L\`u, A.~Kazemnejad, N.~Meade, A.~Patel, D.~Shin, A.~Zambrano, K.~Sta\'nczak, P.~Shaw, et al. AgentRewardBench: Evaluating automatic evaluations of web agent trajectories. arXiv:2504.08942, 2025.
\bibitem{yao2024taubench} S.~Yao, N.~Shinn, P.~Razavi, and K.~Narasimhan. $\tau$-bench: A benchmark for tool-agent-user interaction in real-world domains. arXiv:2406.12045, 2024.
\bibitem{jin2025searchr1} B.~Jin, H.~Zeng, Z.~Yue, J.~Yoon, S.~Arik, D.~Wang, H.~Zamani, and J.~Han. Search-R1: Training LLMs to reason and leverage search engines with reinforcement learning. arXiv:2503.09516, 2025.
\bibitem{yao2023react} S.~Yao, J.~Zhao, D.~Yu, N.~Du, I.~Shafran, K.~Narasimhan, and Y.~Cao. ReAct: Synergizing reasoning and acting in language models. \emph{ICLR}, 2023.
\bibitem{yang2018hotpotqa} Z.~Yang, P.~Qi, S.~Zhang, Y.~Bengio, W.~W.~Cohen, R.~Salakhutdinov, and C.~D.~Manning. HotpotQA: A dataset for diverse, explainable multi-hop question answering. \emph{EMNLP}, 2018.
\end{thebibliography}

\end{document}
